\begin{vidu}
Một vật có khối lượng $m = 10\ kg$ nằm yên trên một mặt sàn ngang. Tác dụng lên vật một lực kéo $\vec{F}$ theo phương ngang. Biết hệ số ma sát nghỉ là $\mu_n = 0,4$, hệ số ma sát trượt là $\mu_t = 0,2$, lấy $g = 10\ m/s^2$.	\begin{listEX}
\item Tính độ lớn lực ma sát nghỉ cực đại.
\item Nếu kéo vật bằng lực $F = 20\ N$, vật có chuyển động không? Tính độ lớn lực ma sát tác dụng lên vật.
\item Nếu kéo vật bằng lực $F = 50\ N$, vật bắt đầu trượt. Tính độ lớn lực ma sát tác dụng và gia tốc chuyển động của vật.
\item Nếu không có lực kéo nữa, vật đang trượt với vận tốc ban đầu $v_0 = 4\ m/s$ thì đi được quãng đường bao nhiêu thì dừng lại?
\end{listEX}

\loigiai{
\begin{itemize}
\item[1.] Lực ma sát nghỉ cực đại:
$$ F_{msn\max} = \mu_n . m . g = 0,4 . 10 . 10 = 40\ N $$

\item[2.] Với $F = 20\ N < F_{msn\max}$ $\Rightarrow$ vật chưa trượt, lực ma sát cân bằng lực kéo:
$$ F_{ms} = 20\ N $$

\item[3.] Với $F = 50\ N > F_{msn\max}$ $\Rightarrow$ vật trượt, lực ma sát là lực ma sát trượt:
$$ F_{mst} = \mu_t . mg = 0,2 . 10 . 10 = 20\ N $$
Gia tốc:
$$ a = \dfrac{F - F_{mst}}{m} = \dfrac{50 - 20}{10} = 3\ m/s^2 $$
\item[4.] Khi không còn lực kéo, vật trượt dưới tác dụng duy nhất của lực ma sát trượt, gây gia tốc:
$$ a = - \mu_t . g = -0,2 . 10 = -2\ m/s^2 $$
Áp dụng công thức chuyển động:
$$ v^2 = v_0^2 + 2as \Rightarrow 0 = 16 + 2 . (-2) . s \Rightarrow s = 4\ m $$
\end{itemize}
}
\end{vidu}

%mac
%\begin{vidu}
%	Một người đi xe đạp có khối lượng tổng cộng $m=86\ kg$ đang chuyển động trên đường nằm ngang với vận tốc $4\ m/s$. Nếu người đi xe ngừng đạp và hãm phanh để giữ không cho các bánh xe quay, xe trượt đi một đoạn đường 2 m thì dừng lại.
%	\begin{itemize}
%		\item[a.] Lực nào đã gây ra gia tốc cho xe? Tính lực này.
%		\item[b.] Tính hệ số ma sát trượt giữa mặt đường và lốp xe. Lấy $g=10\ m/s^2$.
%	\end{itemize}	
%	\haidong{10}\loigiai{
%		Khi tính lực và gia tốc, ta coi người và xe là chất điểm.
%		\begin{itemize}
%			\item[1.] Gia tốc của chuyển động được tính bằng công thức:
%			\begin{align*}
%				a=\dfrac{v^2-v_0^2}{2.s}=\dfrac{0-16}{2.2}=-4\ m/s^2
%			\end{align*}
%			Lực gây ra gia tốc này là lực ma sát trượt của mặt đường tác dụng lên lốp xe:
%			\begin{align*}
%				F=m.a=86.(-4)=-344\ N.
%			\end{align*}
%			\item[2.] Hệ số ma sát trượt giữa lốp xe với mặt đường được tính từ công thức:
%			\begin{align*}
%				F_{ms}=\mu .N\Rightarrow \mu=\dfrac{F_{ms}}{N}
%			\end{align*}
%			Vì ô tô chuyển động trên đường nằm ngang nên: $N=P=m.g\Rightarrow \mu=\dfrac{344}{86.10}=0,4$.
%		\end{itemize}	
%	}	
%\haidong{10} \end{vidu}
%\loai{Tính lực ma sát trượt}
\begin{vidu}%[1P2-7.2K][Loai1]
Biết hệ số ma sát trượt của vật và sàn là $0,4$. Cần kéo một vật trọng lượng 20 N với một lực theo phương ngang có độ lớn bằng bao nhiêu để vật chuyển động đều trên một mặt sàn nằm ngang? 
\choice
{$6\ N$}
{$5\ N$}
{\True $8\ N$}
{$10\ N$}
\haidong{5}\loigiai{
Để vật chuyển động thẳng đều thì $F = F_{ms} = \mu mg = 0,4.20 = 8\ N.$
}
\haidong{10} \end{vidu}
%\loai{Tính lực ma sát nghỉ}
\begin{vidu}%[1P2-7.2K][Loai1]
Tác dụng một lực theo phương ngang có độ lớn $8\ N$ vào một vật có khối lượng $10\ kg$ đang nằm yên trên mặt đất. Biết hệ số ma sát giữa vật và mặt đất là $0,1$. Lấy $g = 10\ m/s^2$, lực ma sát tác dụng vào vật có độ lớn bằng
\choice
{$18\ N$}
{$2\ N$}
{\True $8\ N$}
{$10\ N$}
\haidong{5}\loigiai{
Ta có: $F_{msn\max} = \mu mg = 0,1.10.10 = 10\ N > 8\ N$ nên vật vẫn chưa chuyển động $\Rightarrow$  độ lớn của lực ma sát bằng độ lớn của ngoại lực tác dụng lên vật.
}
\haidong{10} \end{vidu}

%\loai{Tính gia tốc}
\begin{vidu}%[1P2-7.2K][Loai1]
Đẩy một cái thùng có khối lượng $50\ kg $ với lực có độ lớn $150\ N$ theo phương ngang làm thùng chuyển động. Cho biết hệ số ma sát trượt giữa thùng và mặt sàn là $\mu = 0,2$. Lấy $g = 10\ m /s^2 $. Gia tốc của thùng là   
\choice
{\True $a = 1,00\ m/s^2$}
{$a = 1,01\ m/s^2$}
{$a = 1,02\ m/s^2$}
{$a = 1,04\ m/s^2$}
\haidong{4}\loigiai{
Gia tốc chuyển động của vật là: $a = \dfrac{F - F_{ms}}{m} = \dfrac{F - \mu mg}{m} =\dfrac{150 - 0,2.50.10}{50} = 1\ m/s^2$.
}
\haidong{10} \end{vidu}

%\loai{Mặt phẳng ngang}
\begin{vidu}%[1P2-7.2K][Loai1] 
Một vật có khối lượng $m = 8\ kg$ chuyển động thẳng đều trên mặt sàn nằm ngang dưới tác dụng của lực kéo $F = 16\ N$ có phương song song với mặt sàn. Lấy $g=10\ m/s^2$. Xác định lực ma sát và hệ số ma sát giữa vật và mặt sàn. 
\choice
{$F_{ms}=18\ N$ và $\mu=0,2$}
{\True $F_{ms}=16\ N$ và $\mu=0,2$}
{$F_{ms}=16\ N$ và $\mu=0,3$}
{$F_{ms}=18\ N$ và $\mu=0,3$}
\haidong{10}\loigiai{
\immini[thm]{\begin{itemize}
\item Các lực tác dụng lên vật như hình: $\vec{F}$, ${\vec{F}}_{ms}$, $\vec{P}$, $\vec{N}$
\item Chọn hệ trục tọa độ Oxy như hình vẽ
\item Áp dụng định luật II Niutơn:
$\vec{F}+{\vec{F}}_{ms}+\vec{P}+\vec{N}=m\vec{a}$
\item Vì vật chuyển động thẳng đều nên gia tốc a = 0:  $\vec{F}+{\vec{F}}_{ms}+\vec{P}+\vec{N}=0$
\item Chiếu lên Ox ta có: $F-{F}_{ms}=0$$\Rightarrow {F}_{ms}=F=16\ N$
\item Chiếu lên Oy ta có: $N-P=0$$\Rightarrow N=P\Rightarrow {F}_{ms}=\mu P=\mu mg\Rightarrow \mu =\dfrac{{F}_{ms}}{mg}=0,2$. 
\end{itemize}}{\begin{tikzpicture}[scale=1,thick,>=stealth']
\tkzDefPoints{0/1/o}
\node  at  (0,0.32)  [shape=rectangle,draw=red,pattern=dots,minimum height=0.6cm,minimum width=1.4cm, rounded corners=1pt]  {};
\draw(-1.5,0)--(1.5,0);
\draw[very thick,blue,->](0.1,0)--++(90:1.8)node[right,black]{$\vec{N}$};
\draw[fill=white](0.1,0)circle(1.2pt);
\draw[very thick,blue,->](0,0.32)--++(-90:1.8)node[right,black]{$\vec{P}$};
\draw[fill=white](0,0.32)circle(1.2pt);
\draw[very thick,blue,->](-0.7,0)--++(180:1.2)node[above,black]{$\vec{F}_{ms}$};
\draw[fill=white](-0.7,0)circle(1.2pt);
\draw[very thick,blue,->](0.7,0.32)--++(0:1.5)node[above,black]{$\vec{F}$};
\draw[fill=white](0.7,0.32)circle(1.2pt);
\fill[pattern=north east lines](-1.5,0)--(1.5,0)--++(-90:0.1)--++(180:3);
\draw[->](3,0)node[below left]{$O$}--++(90:1.5)node[right]{$y$};
\draw[->](3,0)--++(0:2)node[above]{$x$};
\end{tikzpicture}}

}
\haidong{10} \end{vidu}
%\loai{Tính quãng đường đi được}
\begin{vidu} %[1P2-7.2K][Loai1]
\impicinpar{Một vật khối lượng $m = 400\ g$ đặt trên mặt bàn nằm ngang hình vẽ. Hệ số ma sát trượt giữa vật và mặt bàn là $0,3$. Vật bắt đầu được kéo đi bằng một lực $F = 2\ N$ có phương nằm ngang. Cho $g = 10\ m/s^2$. Quãng đường vật đi được sau $1\ s$ là
\choice
{$0,4\ m$}
{$0,8\ m$}
{\True $1,0\ m$}
{$1,15\ m$}}{\begin{tikzpicture}[draw=Mapcolor, very thick,scale=1,thick,>=stealth']
\tkzDefPoints{0/1/o}
\node  at  (0,0.32)  [shape=rectangle,draw=red,pattern=dots,minimum height=0.6cm,minimum width=1.4cm, rounded corners=1pt]  {};
\draw(-1.5,0)--(1.5,0);
\draw[very thick,blue,->](0.7,0.32)--++(0:1.5)node[above,black]{$\vec{F}$};
\draw[fill=white](0.7,0.32)circle(1.2pt);
\fill[pattern=north east lines](-1.5,0)--(1.5,0)--++(-90:0.1)--++(180:3);
\end{tikzpicture}}
\haidong{10}\loigiai{      
\begin{itemize}
\item Theo định luật II Niutơn ta có: $F - F_{mst}=m.a$ (xét theo phương ngang)
\item Gia tốc của vật: $a = \dfrac{{F - F_{mst} }}{m} = \dfrac{{F - \mu _t mg}}{m} = 2\ m/s^2 $
\item Quãng đường mà vật đi được trong 1 s đầu tiên: $s_1 = \dfrac{{at^2 }}{2}=1\ m.$.
\end{itemize}}
\haidong{10} \end{vidu}
\begin{vidu}%SBT KNTT III.6
An và Bình cùng nhau đẩy một thùng hàng chuyển động thẳng trên sàn nhà. Thùng hàng có khối lượng $120\ kg$. An đẩy với một lực $450\ N$, Bình đẩy với một lực $350\ N$ cùng theo phương ngang. Hệ số ma sát trượt giữa thùng và sàn là $0,2$. Lấy $g=10\ m/s^2$. Gia tốc của thùng là
\choice
{$6,47\ m/s^2$}
{$7,66\ m/s^2$}
{\True $4,67\ m/s^2$}
{$8,94\ m/s^2$}
\haidong{6}\loigiai{
Hợp lực tác dụng lên thùng hàng: $F=F_{\ct{đẩy}}-F_{\ct{ma sát}}$, với $F_{\ct{ma sát}}=\mu mg=240\ N$.
\begin{align*}
F=800-240=560\ N\Rightarrow a=\dfrac{F}{m}=\dfrac{560}{120}=4,67\ m/s^2.
\end{align*}
}
\haidong{10} \end{vidu}
%\newpage
%\loai{Lực theo phương xiên góc}
\begin{vidu} %[1P2-7.2G][Loai1]
Một cái hòm khối lượng $m = 20\ kg$ đặt trên sàn nhà. Người ta kéo hòm bằng một lực hướng chếch lên trên và hợp phương ngang một góc $\alpha=30^\circ$. Hòm chuyển động đều trên sàn nhà nằm ngang. Biết hệ số ma sát trượt giữa hòm và sàn nhà là $0,3$. Cho $g = 9,8\ m/s^2$. Độ lớn của lực kéo hòm là
\choice
{$28,2\ N$}
{\True $57,9\ N$}
{$44,6\ N$}
{$68,5\ N$}
\haidong{13}\loigiai{      
\immini[thm]{\begin{itemize}
\item Chọn hệ trục tọa xOy như hình vẽ
\item Vật chuyển động đều nên ta có: $\vec F + \vec P + \vec N + \vec F_{mst} = \vec 0$
\item Chiếu xuống trục Ox : $F.\cos \alpha - F_{mst} = 0.$
\item Chiếu xuống trục Oy : $F.\sin \alpha - mg + N = 0$.
\item Ngoài ra : $F_{mst} = \mu_t.N$. Thay vào trên ta được: $F.\cos \alpha - \mu_t.N = F.\cos \alpha- \mu_t.(mg - F.\sin \alpha) = 0$ hay $F.(\cos \alpha + \mu_t.\sin \alpha) = \mu_t.mg \Rightarrow$  
$F = \dfrac{{\mu _t mg}}{{\cos \alpha  + \mu _t \sin \alpha }} = 57,87\ N$.
\end{itemize}}{\begin{tikzpicture}[scale=1,thick,>=stealth']
\tkzDefPoints{0/1/o}
\node  at  (0,0.32)  [shape=rectangle,draw=red,pattern=dots,minimum height=0.6cm,minimum width=1.4cm, rounded corners=1pt]  {};
\draw(-1.5,0)--(1.5,0);
\draw[very thick,blue,->](0.1,0)--++(90:1.8)node[right,black]{$\vec{N}$};
\draw[fill=white](0.1,0)circle(1.2pt);
\draw[very thick,blue,->](0,0.32)--++(-90:1.8)node[right,black]{$\vec{P}$};
\draw[fill=white](0,0.32)circle(1.2pt);
\draw[very thick,blue,->](-0.7,0)--++(180:1.2)node[above,black]{$\vec{F}_{ms}$};
\draw[fill=white](-0.7,0)circle(1.2pt);
\draw[very thick,blue,->](0.7,0.32)--++(30:2)node[above,black]{$\vec{F}$};
\draw[fill=white](0.7,0.32)circle(1.2pt);
\fill[pattern=north east lines](-1.5,0)--(1.5,0)--++(-90:0.1)--++(180:3);
\draw[->](3,0)node[below left]{$O$}--++(90:1.5)node[right]{$y$};
\draw[->](3,0)--++(0:2)node[above]{$x$};
\end{tikzpicture}}
}
\haidong{10} \end{vidu}
%\loai{Mặt phẳng nghiêng}

\begin{vidu}%[1P2-7.2G][Loai1]
\immini[thm]{Một vật trượt xuống một mặt phẳng nằm nghiêng dài $5\ m$ và cao $3\ m$. Hệ số ma sát giữa vật và mặt phẳng nghiêng là $0,2$. Lấy $g = 10\ m/s^2$. Gia tốc của vật là
\haicot
{$5\ m/s^2$}
{$3,5\ m/s^2$}
{\True $4,4\ m/s^2$}
{$3,9\ m/s^2$}}{\begin{tikzpicture}[scale=1,thick,>=stealth']
\tkzInit[ymin=-4,ymax=0.5,xmin=-0.5,xmax=5.5]\tkzClip
\begin{scope}[shift={(0,0)},rotate=-30]
\node  at  (0.3,0)  [shape=rectangle,draw=red,pattern=dots,minimum height=0.5cm,minimum width=0.6cm, rounded corners=1pt,rotate=-30]  {};
\draw(0,-0.27)coordinate(a)--++(0:6)coordinate(b)--++(210:5.2)coordinate(c);
\draw(0,-0.27)--++(-60:3);
\fill[pattern=north east lines](b)--(c)--++(300:0.1)--++(390:5.2);
\tkzMarkAngles[size=0.5cm](a,b,c)
\tkzLabelAngle[pos=-0.8](a,b,c){$\alpha$}
\end{scope}
\end{tikzpicture}}
\haidong{10}\loigiai{
\immini[thm]{\begin{itemize}
\item Ta có: $F_{ms} = \mu mg cos \alpha ;\  F_P = mg sin \alpha$ 
\item Gia tốc chuyển động của vật là:\\
$a = \dfrac{F_P - F_{ms}}{m} = \dfrac{mg \sin \alpha  - \mu mg \cos \alpha}{m} = g(\sin \alpha  - \mu \cos \alpha)\\ 
= 10. \left(\dfrac{3}{5} - 0,2.\dfrac{\sqrt {5^2 - 3^2 }}{5}\right)  = 4,4\ m/s^2$.
\end{itemize}}{\begin{tikzpicture}[scale=1,thick,>=stealth']
\tkzInit[ymin=-4,ymax=0.5,xmin=-0.5,xmax=5.5]\tkzClip
\begin{scope}[shift={(0,0)},rotate=-30]
\node  at  (2.3,0)  [shape=rectangle,draw=red,pattern=dots,minimum height=0.5cm,minimum width=0.6cm, rounded corners=1pt,rotate=-30]  {};
\draw[very thick,blue,->](2,-0.25)coordinate(o)--++(180:0.5)node[below,black,xshift=-10pt]{$\vec{F}_{ms}$};
\draw[fill=white](2,-0.25)circle(1.2pt);
\draw[very thick,blue,->](2.3,0)--++(0:0.7)node[above,black,xshift=10pt]{$\vec{F}_p$}coordinate(px);
\draw[very thick,blue,->](2.3,0)coordinate(o')--++(90:1.2)node[right,black]{$\vec{R}$};
\draw[very thick,blue,->](2.3,0)--++(-90:1.2)node[left,black]{$\vec{N}$}coordinate(n);
\draw[very thick,blue,->](2.3,0)--++(-60:1.4)node[right,black,yshift=5pt]{$\vec{P}$}coordinate(p);
\draw[dashed,thin](px)--(p)--(n);
\draw[fill=white](2.3,0)circle(1.2pt);
\draw(0,-0.27)coordinate(a)--++(0:6)coordinate(b)--++(210:5)coordinate(c);
\fill[pattern=north east lines](b)--(c)--++(300:0.1)--++(390:5);
\tkzMarkAngles[size=0.5cm](a,b,c)
\tkzLabelAngle[pos=-0.8](a,b,c){$\alpha$}
\tkzMarkAngles[size=0.5cm](px,p,o')
\tkzLabelAngle[pos=0.7](px,p,o'){$\alpha$}
\end{scope}
\end{tikzpicture}}
}
\haidong{10} \end{vidu}
%%\loai{Lực căng}
%\begin{vidu}%[1P2-7.2G][Loai1]
%\immini[thm]{Một vật được bố trí như hình vẽ, vật có khối lượng $m = 500$ g, $\alpha  = 45^\circ$, dây AB song song với mặt phẳng nghiêng, hệ số ma sát nghỉ giữa vật và mặt phẳng nghiêng là $\mu_n = 0,5$. Cho $g = 9,8\ m/s^2$. Độ lớn của lực ma sát giữa vật và mặt phẳng nghiêng và lực căng của dây lần lượt là
%\choice
%{3,46 N; 3,46 N}
%{\True 1,73 N; 1,73 N}
%{1,8 N; 1,8 N}
%{1,73 N; 3,46 N}}{\begin{tikzpicture}[scale=1,thick,>=stealth']\\
%\tkzInit[ymin=-3.5,ymax=0.5,xmin=-0.5,xmax=5.5]\tkzClip
%\begin{scope}[shift={(0,0)},rotate=-30]
%\node  at  (3.3,0)  [shape=rectangle,draw=red,pattern=dots,minimum height=0.5cm,minimum width=0.6cm, rounded corners=1pt,label=above:$m$,rotate=-30]  {};
%\draw(0,0) -- (3,0);
%\draw(0,0.27)--(0,-0.27)coordinate(a)--++(0:6)coordinate(b)--++(210:5)coordinate(c);
%\fill[pattern=horizontal lines](0,0.27)--(0,-0.27)--++(180:0.1)--++(90:0.54);
%\fill[pattern=north east lines](b)--(c)--++(300:0.1)--++(390:5);
%\tkzMarkAngles[size=0.5cm](a,b,c)
%\tkzLabelAngle[pos=-0.8](a,b,c){$\alpha$}
%\end{scope}
%\end{tikzpicture}}
%\haidong{10}\loigiai{
%\immini[thm]{\begin{itemize}
%\item Ta biểu diễn các lực tác dụng lên vật như hình.
%\item Ta có:  $F_{P} = mg\sin \alpha  = 0,5.9,8.\sin 45^\circ = 3,46\ N$.
%\item $F_{msn\max}  = \mu_nmg\cos \alpha  = 0,5.0,5.9,8.\cos 45^\circ = 1,73 \ N$.
%\item $F_{P_1}$ có xu hướng kéo vật trượt xuống, giá trị của nó lớn hơn giá trị lớn nhất của lực ma sát nghỉ $\Rightarrow$ lực ma sát nghỉ đạt tới giá trị cực đại bằng $1,73\ N$.
%\item Vật cân bằng $\Rightarrow$ $T = {F_P} - F_{ms} = 3,46 - 1,73 = 1,73\ N$.
%\end{itemize}}{\begin{tikzpicture}[scale=1,thick,>=stealth']
%\tkzInit[ymin=-3.5,ymax=0.5,xmin=-0.5,xmax=5.5]\tkzClip
%\begin{scope}[shift={(0,0)},rotate=-30]
%\node  at  (3.3,0)  [shape=rectangle,draw=red,pattern=dots,minimum height=0.5cm,minimum width=0.6cm, rounded corners=1pt,rotate=-30]  {};
%\draw(0,0) -- (3,0);
%\draw[very thick,blue,->](3,0)coordinate(o)--++(180:0.5)node[above,black,xshift=10pt]{$\vec{T}$};
%\draw[fill=white](3,0)circle(1.2pt);
%\draw[very thick,blue,->](3,-0.25)--++(180:0.4)node[below,black,xshift=-10pt]{$\vec{F}_{ms}$};
%\draw[fill=white](3,-0.25)circle(1.2pt);
%\draw[very thick,blue,->](3.3,0)--++(0:0.7)node[above,black,xshift=10pt]{$\vec{F}_{P1}$}coordinate(px);
%\draw[very thick,blue,->](3.3,0)--++(90:1.2)node[right,black]{$\vec{R}$};
%\draw[very thick,blue,->](3.3,0)--++(-90:1.2)node[left,black]{$\vec{N}$}coordinate(n);
%\draw[very thick,blue,->](3.3,0)--++(-60:1.4)node[right,black,yshift=5pt]{$\vec{P}$}coordinate(p);
%\draw[dashed,thin](px)--(p)--(n);
%\draw[->](o)--++(0:2)node[above right,black]{$x$};
%\draw[fill=white](3.3,0)circle(1.2pt);
%\draw(0,0.27)--(0,-0.27)coordinate(a)--++(0:6)coordinate(b)--++(210:5)coordinate(c);
%\fill[pattern=horizontal  lines](0,0.27)--(0,-0.27)--++(180:0.1)--++(90:0.54);
%\fill[pattern=north east lines](b)--(c)--++(300:0.1)--++(390:5);
%\tkzMarkAngles[size=0.5cm](a,b,c)
%\tkzMarkAngles[size=0.5cm](px,p,o)
%\tkzLabelAngle[pos=0.7](px,p,o){$\alpha$}
%\end{scope}
%\end{tikzpicture}}
%}
%\haidong{10} \end{vidu}